\scanpage{6}
% Two short horizontal strokes at the upper right.
\margin{\mbox{-- --}}

\deleted{\unsure{catégories|catégorie}}\inserted{champs} de Picard \unsure{épinglés|épinglées} par $M.N$, $\sigma$ étant l'invariant qui \inserted{s'obtient en\unsure{,}} associant à toute section \inserted{$L$} d'un champ de Picard la symétrie de $\deleted{\unsure{N|L}}L\otimes L$, \unsure{interprétée} comme section de ${}_{2}N$. Je sais prouver (sauf erreur) que tout hom. $M\to{}_{2}N$ provient d'un champ de Picard convenable (épinglé par $M,N$) (a priori l'obstruction est dans $\mathrm{Ext}^{3}(X;M,N)$, mais un argument ``universel'' prouve qu'elle est nulle). Cela prouve que l'extension $(*)$ est bien proche d'être splittée : \deleted{\unsure{un|la}} tte section du troisième faisceau, sur un objet quelconque de $X$, se \unsure{remonte} -- en d'autres termes l'extension a une section ``ensembliste''. Bien sûr, il y a mieux en fait : toute section sur un $U\in\mathrm{Ob}\,X$ ``provient'' d'un élément de ${\unsure{\mathscr{H}|\mathcal{H}}}^{2}(U,{\unsure{E|\underline{F}}}'(M\unsure{.}\;N\unsure{)|;|,}))$ (hypercohomologie $\mathscr{H}^{2}$).

% Reproduction boundary strokes are visible at the left and below the text; a short thick vertical reproduction mark is at the lower right.
% archive sheet number: 63
